\roteiro{3}

% Revelação em TikZ sem deslocar nada: o elemento existe em todas as páginas do frame,
% transparente fora do intervalo <...> (estilo usual do beamer com TikZ).
\tikzset{
  invisivel/.style={opacity=0, text opacity=0},
  visivel em/.style={alt={#1{}{invisivel}}},
  alt/.code args={<#1>#2#3}{\alt<#1>{\pgfkeysalso{#2}}{\pgfkeysalso{#3}}}}

% Slide 8: visão geral (relatório §4.1, Fig. 7)
% A Fig. 7 (images/figure7-overview.png, 1376 x 1380 px) rende texto de ~3,6 pt na altura
% útil do slide; o desenho abaixo repete as suas caixas e setas em TikZ, com texto legível.
\begin{frame}{Visão geral}
  \centering
  \begin{tikzpicture}[x=1mm, y=1mm, font=\footnotesize,
      caixa/.style={draw=ufscborda, fill=white, rounded corners=0.6mm, inner sep=2pt,
                    align=center, text=ufsctexto},
      rotulo/.style={font=\footnotesize\bfseries, text=ufscazul},
      seta/.style={->, >=stealth, draw=ufscazul!70, thick}]
    \node[rotulo] at (20, 31) {Entrada};
    \node[rotulo, visivel em=<2->] at (68, 31) {Processamento};
    \node[rotulo, visivel em=<3->] at (120, 31) {Saída};
    % entrada
    \node[caixa, fill=ufscfundo, text width=40mm] (imp) at (20, 18)
      {\textbf{arquivo de importação}\\[-1pt]{\scriptsize\texttt{import\_config.json}}\\[-1pt]
       {\scriptsize fontes e mapeamento}};
    \node[caixa, fill=ufscfundo, text width=40mm] (con) at (20, 2)
      {\textbf{arquivo de conexão}\\[-1pt]{\scriptsize\texttt{dbms\_config.json}}\\[-1pt]
       {\scriptsize SGBDs e como alcançá-los}};
    \node[caixa, fill=ufscfundo, text width=40mm] (csv) at (20, -14)
      {\textbf{arquivos de dados}\\[-1pt]{\scriptsize um ou mais CSV}};
    % processamento
    \begin{scope}[visivel em=<2->]
      \node[caixa, fill=ufscazul, draw=ufscazul, text=white, text width=38mm,
            minimum height=16mm] (tool) at (68, 2)
        {\bfseries Polyglot Import CSV\\[1pt]{\scriptsize\mdseries comando CLI:\\
         valida e importa}};
      \draw[seta] (imp.east) -- (tool.west);
      \draw[seta] (con.east) -- (tool.west);
      \draw[seta] (csv.east) -- (tool.west);
    \end{scope}
    % saída
    \begin{scope}[visivel em=<3->]
      \foreach \sgbd/\tipo [count=\i] in {PostgreSQL/relacional, MongoDB/documento,
          Cassandra/colunar, Redis/chave-valor, Neo4j/grafo} {
        \node[caixa, text width=34mm, minimum height=5.6mm] (s\i) at (120, {26-8*\i})
          {\textbf{\sgbd}\enspace{\scriptsize\color{black!60}\tipo}};
        \draw[seta] (tool.east) -- (s\i.west);
      }
    \end{scope}
  \end{tikzpicture}\par
  \medskip
  \onslide<3->{\footnotesize O comando valida as entradas e só então importa para os SGBDs
   presentes na configuração.\par}
\end{frame}

% Slide 9: arquivo de importação (relatório §4.2.1; trecho de data/ecommerce/import_config.json)
% O trecho mantém as chaves e o aninhamento do arquivo real; "..." marca o que foi omitido
% (entidades categories e products antes de inventory, relationships, os outros SGBDs).
\begin{frame}[fragile]{Arquivo de importação}
  \destaque{Arquivo de importação}: JSON com as fontes (CSVs) e o que gravar em cada SGBD.\par
  \smallskip
  \begingroup\footnotesize Trecho de \texttt{import\_config.json}: a fonte \texttt{stock}
  alimenta a tabela \texttt{inventory}.\par\endgroup
  \smallskip
  \begin{minipage}[t]{0.66\textwidth}
\begin{lstlisting}[style=jsonslide, aboveskip=0pt, belowskip=0pt]
"sources": {
  "stock": "ecommerce_stock.csv",
  "purchase": "ecommerce_purchase.csv", ...
},
"postgres": {
  "entities": { ...
    "inventory": {
      "source": "stock",
      "columns": {
        "product_id": { "is_key": true, "db_type": "BIGINT" },
        "quantity_available": { "db_type": "BIGINT" }, ...
      } }, ...
  }, ...
}, ...
\end{lstlisting}
  \end{minipage}\hfill
  \begin{minipage}[t]{0.32\textwidth}
    \footnotesize
    \begin{itemize}
      \item \texttt{sources}: dá um nome a cada CSV.
      \item \texttt{postgres}: um bloco por SGBD de destino.
      \item \texttt{entities}: tabelas, coleções, nós\ldots
      \item \texttt{columns}: colunas, com chave (\texttt{is\_key}) e tipo
            (\texttt{db\_type}).
    \end{itemize}
  \end{minipage}
\end{frame}

% Slide 10: conexão e validação prévia (relatório §4.2.3, §4.2.4, Fig. 10)
% A Fig. 10 (images/figure10-config-pipeline.png) é um retrato alto; o fluxo foi redesenhado
% na horizontal, com as três camadas de validação reveladas uma a uma.
\begin{frame}{Conexão e validação prévia}
  \destaque{Arquivo de conexão}: JSON que diz onde alcançar cada SGBD.\par
  \smallskip
  Antes de gravar, a validação passa por \destaque{três camadas}:\par
  \medskip
  \centering
  \begin{tikzpicture}[x=1mm, y=1mm, font=\scriptsize,
      caixa/.style={draw=ufscborda, fill=white, rounded corners=0.6mm, inner sep=2pt,
                    align=center, text=ufsctexto, minimum height=11mm},
      entrada/.style={caixa, fill=ufscfundo, text width=19mm},
      camada/.style={caixa, text width=25mm},
      seta/.style={->, >=stealth, draw=ufscazul!70, thick},
      falha/.style={->, >=stealth, draw=red!55!black, thick}]
    % entradas
    \node[entrada] (imp) at (10.5, 7) {arquivo de importação};
    \node[entrada] (con) at (10.5, -7) {arquivo de conexão};
    % camada 1
    \begin{scope}[visivel em=<2->]
      \node[camada] (c1) at (39, 0)
        {\textbf{1. JSON Schema}\\ estrutura e tipos de cada arquivo};
      \draw[seta] (imp.east) -- (c1.west);
      \draw[seta] (con.east) -- (c1.west);
      \node[caixa, draw=red!55!black, fill=red!6, text width=98mm, minimum height=7mm]
        (erro) at (70.5, -19)
        {\textbf{Falha de validação:} a execução termina e \textbf{nada é gravado} nos
         SGBDs.};
      \draw[falha] (c1.south) -- (c1.south |- erro.north);
    \end{scope}
    % camada 2
    \begin{scope}[visivel em=<3->]
      \node[camada] (c2) at (70.5, 0)
        {\textbf{2. Consistência entre arquivos}\\ todo SGBD da importação está na
         conexão?};
      \draw[seta] (c1.east) -- (c2.west);
      \draw[falha] (c2.south) -- (c2.south |- erro.north);
    \end{scope}
    % camada 3
    \begin{scope}[visivel em=<4->]
      \node[entrada, minimum height=7mm] (csv) at (102, 19) {CSVs (fontes)};
      \node[camada] (c3) at (102, 0)
        {\textbf{3. Validação cruzada com os dados}\\ colunas, chaves e arestas existem?};
      \draw[seta] (c2.east) -- (c3.west);
      \draw[seta] (csv.south) -- (c3.north);
      \draw[falha] (c3.south) -- (c3.south |- erro.north);
    \end{scope}
    % importação
    \begin{scope}[visivel em=<5->]
      \node[caixa, fill=ufscazul, draw=ufscazul, text=white, text width=17mm] (ok)
        at (129.5, 0) {\textbf{importação} por SGBD};
      \draw[seta] (c3.east) -- (ok.west);
    \end{scope}
  \end{tikzpicture}
\end{frame}

% Slide 11: cenário e-commerce (relatório §4.2, tab:eventos-ecommerce; §4.2.1, modos;
% §4.3.2, o exemplo usa o modo multifonte). Os blocos sources vêm de import_config.json e
% import_config_combined.json (iguais a lst:sources-multi e lst:sources-combined).
\begin{frame}[fragile]{Cenário e-commerce}
  \begingroup\footnotesize\bfseries\color{ufscazul} Tipos de evento do cenário de
  e-commerce\hfill\scriptsize\mdseries\color{black!60} todo evento registra o instante e o
  usuário envolvido\par\endgroup
  \smallskip
  \footnotesize
  \setlength{\tabcolsep}{4pt}
  \setlength{\aboverulesep}{0.15ex}\setlength{\belowrulesep}{0.3ex}
  \begin{tabular}{@{}>{\ttfamily}l l E{0.55\textwidth}@{}}
    \toprule
    \normalfont\bfseries Tipo & \bfseries Evento & \bfseries Demais dados registrados \\
    \midrule
    stock           & reposição de estoque & produto, categoria, quantidade disponível e preço \\
    purchase        & compra               & pedido, produto, contraparte da venda, pagamento e avaliação \\
    select\_product & seleção de produto   & produto selecionado \\
    add\_to\_cart   & adição ao carrinho   & carrinho, produto e quantidade \\
    \bottomrule
  \end{tabular}\par
  \smallskip
  \onslide<2->
  \begingroup\tcbset{equal height group=modos}
  \begin{minipage}[t]{0.49\textwidth}
    \begin{painel}{Multifonte\enspace{\mdseries\color{black!60}um CSV por tipo de evento}}{}
\begin{lstlisting}[style=jsonslide, aboveskip=0pt, belowskip=0pt]
"sources": {
  "stock": "ecommerce_stock.csv",
  "purchase": "ecommerce_purchase.csv", ...
}
\end{lstlisting}
    \end{painel}
  \end{minipage}\hfill
  \begin{minipage}[t]{0.49\textwidth}
    \begin{painel}{Combinado\enspace{\mdseries\color{black!60}um CSV com coluna de
      origem}}{}
\begin{lstlisting}[style=jsonslide, aboveskip=0pt, belowskip=0pt]
"sources": {
  "ecommerce": {
    "file": "ecommerce_join.csv",
    "origin_column": true } }
\end{lstlisting}
    \end{painel}
  \end{minipage}\par
  \endgroup
  \smallskip
  Nos dois modos, o resto do arquivo de importação é igual. O exemplo usa o multifonte.\par
\end{frame}

% Slide 12: o mesmo cenário em cinco modelos de dados (relatório §4.3.3)
% Valores e nomes de campo copiados de tab:pg-inventory, lst:mongo-doc,
% tab:cassandra-particao, lst:redis-pares e lst:neo4j-grafo; "..." marca o que foi omitido.
% A linha de stock do Product2 não alimenta o Redis (tab:exemplo-ecommerce): cada painel
% indica a fonte de onde veio. Corpo dos painéis em Helvetica Narrow: nomes como
% quantity_available não cabem em mono \scriptsize num quinto da largura.
\newcommand{\estreita}{\fontfamily{phv}\fontseries{mc}\selectfont}
\newcommand{\pfonte}[1]{{\color{black!60}$\leftarrow$ #1}\par}
\newcommand{\pent}[1]{\textbf{#1}\par}
\newcommand{\pind}{\hspace*{0.4em}}
\newcommand{\prot}[1]{{\color{black!60}#1}}
% Corpo em \fontsize{8}{9}: mesmo corpo do \scriptsize, entrelinha menor (como no slide 7).
\newcommand{\modelo}[4]{% #1 passo, #2 SGBD, #3 modelo de dados, #4 corpo
  \onslide<#1->{\begin{minipage}[t]{0.194\textwidth}
    \footnotesize
    \begin{painel}{#2}{#3}\fontsize{8}{9}\selectfont\estreita\shorthandoff{"}#4\end{painel}
  \end{minipage}}}
\begin{frame}{O mesmo cenário em cinco modelos de dados}
  \begin{tcolorbox}[colback=ufscfundo, colframe=ufscazul, sharp corners, boxrule=0pt,
      leftrule=1mm, left=2mm, right=2mm, top=1mm, bottom=1mm, boxsep=0pt]
    \footnotesize
    \destaque{Fio condutor:} \texttt{user1} (\texttt{user\_id} \texttt{7fd7...}) e
    \texttt{Product2} (\texttt{product\_id} 2)\hfill
    {\scriptsize\color{black!60}$\leftarrow$ indica a fonte (CSV)}\par
  \end{tcolorbox}
  \smallskip
  \begingroup\tcbset{equal height group=modelos}
  \modelo{1}{PostgreSQL}{relacional}{%
    \pfonte{stock}
    \pent{inventory}
    \begin{tabular*}{\linewidth}{@{}l@{\extracolsep{\fill}}r@{}}
      product\_id & 2\\
      quantity\_available & 2\\
      \multicolumn{2}{@{}l@{}}{last\_restock\_date}\\
      \multicolumn{2}{@{}r@{}}{02/11/2023 03:30}\\
      price & 50.5
    \end{tabular*}}\hfill
  \modelo{2}{MongoDB}{documento}{%
    \pfonte{stock}
    \pent{product\_catalog}
    \begin{tabular}{@{}l@{}}
      \{"product\_id": "2",\\
      \pind ...,\\
      \pind "category": \{...\},\\
      \pind "stock": \{\\
      \pind\pind "quantity\_available":\\
      \pind\pind\pind "2", ...\}\}
    \end{tabular}}\hfill
  \modelo{3}{Cassandra}{colunar}{%
    \pfonte{as quatro fontes}
    \pent{user\_activity\_log}
    \prot{partição do user1}\par
    user\_id = 7fd7...\par
    \begin{tabular}{@{}l@{\hspace{3pt}}l@{}}
      \multicolumn{2}{@{}l@{}}{\prot{event\_time}}\\
      & \prot{event\_type}\\
      \multicolumn{2}{@{}l@{}}{\textbf{02/11}}\\
      03:30 & stock\\
      \multicolumn{2}{@{}l@{}}{\textbf{15/11}}\\
      15:09 & select\_product\\
      15:17 & add\_to\_cart\\
      15:46 & purchase
    \end{tabular}}\hfill
  \modelo{4}{Redis}{chave-valor}{%
    \pfonte{add\_to\_cart}
    \pent{shopping\_cart}
    \begin{tabular}{@{}l@{}}
      \prot{chave} user:7fd7...:cart\\
      \prot{valor} \{...,\\
      "cart\_product\_id": "3",\\
      "cart\_quantity": "1"\}
    \end{tabular}\par
    \smallskip
    \pfonte{select\_product}
    \pent{user\_session}
    \begin{tabular}{@{}l@{}}
      \prot{chave} 7fd7...\\
      \prot{valor} \{"user\_name":\\
      \pind "user1", ...\}
    \end{tabular}}\hfill
  \modelo{5}{Neo4j}{grafo}{%
    \pfonte{stock e purchase}
    \begin{tabular}{@{}l@{}}
      (:User \{...,\\
      \pind user\_name: "user1",\\
      \pind ...\})\\
      -[:PURCHASED \{\\
      \pind order\_number:\\
      \pind\pind "ORD1", ...\}]->\\
      (:Product \{\\
      \pind product\_id: "2", ...\})
    \end{tabular}}\par
  \endgroup
\end{frame}
